\bibitem{ResidueIntake} Received AI-generated web-session calculation,
\emph{Fable arithmetic residues and full-Weil rational tests},
20 September 2026. The complete unchanged incoming source is retained as
\texttt{source\_intake/INCOMING.md} in this edition's source archive.
No author name was supplied. Equations (1)--(17) and (33)--(36)
are the specific inputs independently calculated here; AR13 corrects
the omitted completion-factor modulus after incoming equation (16).
The separately developed full-Weil coefficient and tail calculation
is not assumed here.
\bibitem{DLMFXi} T. M. Apostol, \emph{Zeta and Related Functions},
NIST Digital Library of Mathematical Functions, version 1.2.8,
15 September 2026, \url{https://dlmf.nist.gov/25.4}.
Original equation TeX 25.4.3 and 25.4.4 supplies the reflection
identity and exact Riemann completion; 25.2.12 records the canonical
product. The unchanged equation sources are included with their hashes.
The source notes for 25.4 refer to Apostol (1976), p.~260.
The classical line-one argument associated with J. Hadamard and
C.-J. de la Vall\'ee Poussin (1896) is reproduced in AX4--6;
it is not claimed as a new general zeta theorem.
\bibitem{DLMFTheta} W. P. Reinhardt and P. L. Walker,
\emph{Theta Functions}, NIST Digital Library of Mathematical Functions,
version 1.2.8, 15 September 2026,
\url{https://dlmf.nist.gov/20.2.E3} and
\url{https://dlmf.nist.gov/20.7.E32}.
The original unchanged equation TeX supplies the theta series and
modular transformation. AX5 derives the entire theta integral and
the stated numerical bound from these two formulas.
\bibitem{DLMFGamma} R. A. Askey and R. Roy,
\emph{Gamma Function}, NIST Digital Library of Mathematical Functions,
version 1.2.8, 15 September 2026,
\url{https://dlmf.nist.gov/5.4.E4}.
Original equation TeX 5.4.4 is retained unchanged. Its exact
Gamma modulus supplies the equality and exponential factor in AX7.
